\section{Peripheral spectrum and primitivity}

\begin{definition}[Peripheral spectrum]\label{def:peripheral_spectrum}
    \lean{peripheralSpectrum}
    \leanok
    The \emph{peripheral spectrum} of a bounded linear operator $T$
    is the set of spectral values $\mu$ with
    $|\mu|$ equal to the spectral radius of $T$.
\end{definition}

\begin{definition}[Peripheral eigenvalues]\label{def:peripheral_eigenvalues}
    \lean{peripheralEigenvalues}
    \leanok
    The \emph{peripheral eigenvalues} of a linear map $E$ on
    a finite-dimensional space
    are the eigenvalues $\lambda$ on the unit circle:
    $|\lambda| = 1$.
    For a channel, the spectral radius is~$1$
    \cite[Proposition~6.1]{Wolf2012Quantum}, so these coincide with
    the eigenvalues of maximal modulus.

    The condition $|\lambda| = 1$ is appropriate for channels, since the
    spectral radius of a channel is~$1$; for a general linear map with
    spectral radius $r \neq 1$ the condition would generalise to
    $|\lambda| = r$.
\end{definition}

\begin{lemma}[Unit spectral radius forces positive matrix size]
    \label{lem:spectral_radius_one_matrix_dim_pos}
    \lean{matrix_dim_ne_zero_of_spectralRadius_eq_one}
    \leanok
    Let $T:\MN D\to\MN D$ be a bounded linear map. If its spectral radius is
    one, then $D\geq1$.
\end{lemma}
\begin{proof}\leanok
    If $D=0$, then the matrix space and its endomorphism algebra contain only
    zero, so every endomorphism has spectral radius zero.
\end{proof}

\begin{remark}[Peripheral conventions for channels]%
    \label{rem:peripheral_conventions_channels}
    \leanok
    \uses{def:peripheral_spectrum, def:peripheral_eigenvalues}
    The spectral-radius convention of
    Definition~\ref{def:peripheral_spectrum} is the usual bounded-operator
    notion. For channels the spectral radius is~$1$, so the channel arguments
    below use the unit-circle set of Definition~\ref{def:peripheral_eigenvalues}.
\end{remark}

\begin{definition}[Primitive map]\label{def:primitive_channel}
    \lean{IsPrimitive}
    \leanok
    \uses{def:peripheral_eigenvalues}
    A linear map is \emph{primitive} if its only peripheral eigenvalue
    is $\lambda = 1$, i.e.\
    $\mathrm{peripheral}(E) = \{1\}$.
    For channels this corresponds to \cite[Theorem~6.7]{Wolf2012Quantum}.
    This definition applies to any linear endomorphism, not only
    to channels.
\end{definition}

\section{Fixed-point projection}

\begin{definition}[Fixed-point projection]\label{def:fixed_point_proj}
    \lean{fixedPointProj}
    \leanok
    Given a linear map $E:\MN{D}\to\MN{D}$ and a fixed point
    $\rho$ with $E(\rho)=\rho$ and $\tr(\rho)\neq0$,
    the \emph{fixed-point projection} is the rank-one map
    \begin{align}
        P(X)
         & =\frac{\tr(X)}{\tr(\rho)}\rho.
        \label{eq:channel_fixed_point_proj}
    \end{align}
    This projects onto the span of $\rho$ along the kernel of the
    trace functional. When the fixed point is unique up to scaling,
    this span is the full fixed-point space.
\end{definition}

\begin{theorem}[Power decomposition]\label{thm:pow_decomp}
    \lean{pow_eq_fixedPointProj_add_compl_pow}
    \leanok
    \uses{def:fixed_point_proj, def:trace_preserving}
    If $E$ is trace-preserving and $E(\rho)=\rho$, then, for $n\geq1$,
    \begin{align}
        E^n
         & =P+(E-P)^n,
        \label{eq:channel_power_decomp}
    \end{align}
    where $P$ is the fixed-point projection.
\end{theorem}

\begin{proof}
    \leanok
    \uses{def:fixed_point_proj}
    By (\ref{eq:channel_fixed_point_proj}), $P$ is idempotent.
    The fixed-point relation $E(\rho)=\rho$ gives $E\circ P=P$, and trace
    preservation gives $P\circ E=P$. These imply
    $(E-P)\circ P=0$ and $P\circ(E-P)=0$, so in the binomial expansion of
    $(P+(E-P))^n$ all cross terms vanish, yielding
    (\ref{eq:channel_power_decomp}).
\end{proof}

\section{Peripheral eigenvalues and powering}


\begin{theorem}[$1$ is a peripheral eigenvalue]\label{thm:one_peripheral}
    \lean{one_mem_peripheralEigenvalues}
    \leanok
    \uses{def:peripheral_eigenvalues}
    If $E$ has a nonzero fixed point $\rho\neq0$ with $E(\rho)=\rho$,
    then $1$ is a peripheral eigenvalue of $E$.
\end{theorem}

\begin{proof}\leanok
    $\rho$ is an eigenvector with eigenvalue~$1$, and $|1|=1$.
\end{proof}

\begin{theorem}[Finiteness of peripheral eigenvalues]\label{thm:peripheral_finite}
    \lean{peripheralEigenvalues_finite}
    \leanok
    \uses{def:peripheral_eigenvalues}
    On a finite-dimensional space, the set of peripheral eigenvalues
    is finite.
\end{theorem}

\begin{proof}\leanok
    Peripheral eigenvalues are a subset of all eigenvalues,
    which is a finite set in finite dimensions.
\end{proof}

\begin{theorem}[Power-stable peripheral eigenvalues are roots of unity]\label{thm:peripheral_roots_of_unity}
    \lean{peripheral_isRootOfUnity_of_pow_eigenvalue}
    \leanok
    \uses{def:peripheral_eigenvalues}
    Let $E$ be a linear endomorphism on a finite-dimensional space.
    If $\mu$ is a peripheral eigenvalue of $E$ and $\mu^n$ is an eigenvalue
    of $E$ for every $n\geq1$, then $\mu$ is a root of unity.
\end{theorem}

\begin{proof}
    \leanok
    Since $\mu^n$ is an eigenvalue for every positive~$n$,
    the pigeonhole principle on the finite eigenvalue set gives
    $\mu^a=\mu^b$ for some $a\neq b$, hence $\mu^{|a-b|}=1$.
\end{proof}

\begin{theorem}[Peripheral powers imply root of unity]\label{thm:peripheral_powers}
    \lean{peripheral_isRootOfUnity_of_closed_powers}
    \leanok
    \uses{def:peripheral_eigenvalues}
    If the set of peripheral eigenvalues of a linear endomorphism~$E$ on a
    finite-dimensional space is \emph{closed under powers}
    ($\mu\in\mathrm{peripheral}(E)$ and $n\geq1$ imply
    $\mu^n\in\mathrm{peripheral}(E)$), then every peripheral eigenvalue is a
    root of unity.
\end{theorem}

\begin{proof}
    \leanok
    \uses{thm:peripheral_roots_of_unity}
    The closure hypothesis provides
    $\mu^n\in\mathrm{peripheral}(E)$ for all positive~$n$,
    which in particular means $\mu^n$ is an eigenvalue.
    Theorem~\ref{thm:peripheral_roots_of_unity} then gives
    $\mu^p=1$ for some $p>0$.
\end{proof}

\subsection{Periodicity removal by powering}\label{subsec:periodicity_removal_power}

\begin{remark}\label{rem:periodicity_removal_spectral}\leanok
    The results in this subsection are purely spectral: they use only
    finiteness of a set of complex numbers and the spectral mapping theorem
    for powers. In particular, they do not rely on complete positivity.
    For transfer maps, replacing a tensor by its $p$-blocked tensor replaces
    $\E_A$ by $\E_A^p$, so this powering step is the operator-theoretic
    form of blocking.
\end{remark}

\begin{lemma}[Common exponent for a finite set of roots of unity]\label{lem:common_power_eq_one}
    \lean{exists_common_power_eq_one_of_finite}
    \leanok
    Let $s\subseteq\C$ be a finite set.
    Assume that for every $\mu\in s$ there exists an exponent $p_\mu>0$
    with $\mu^{p_\mu}=1$.
    Then there exists $p>0$ such that $\mu^p=1$ for all $\mu\in s$.
\end{lemma}

\begin{proof}\leanok
    Take $p$ to be the least common multiple of the exponents $p_\mu$.
\end{proof}

\begin{lemma}[Powering collapses peripheral eigenvalues]\label{lem:peripheral_pow_singleton}
    \lean{peripheralEigenvalues_pow_eq_singleton}
    \leanok
    \uses{def:peripheral_eigenvalues}
    Let $E:\MN{D}\to\MN{D}$ be a linear map, and assume that $E$ has a
    nonzero fixed point $\rho\neq0$. Let $p>0$.
    If every peripheral eigenvalue $\mu$ of $E$ satisfies $\mu^p=1$,
    then $\mathrm{peripheral}(E^p)=\{1\}$.
\end{lemma}

\begin{proof}\leanok
    \uses{thm:one_peripheral}
    The fixed point $\rho$ gives $1\in\mathrm{peripheral}(E^p)$.
    For the reverse inclusion, let $\nu\in\mathrm{peripheral}(E^p)$.
    Spectral mapping gives $\nu=\mu^p$ for some $\mu\in\spec(E)$.
    From $|\nu|=1$ and $p>0$ we obtain $|\mu|=1$, hence
    $\mu\in\mathrm{peripheral}(E)$.
    The hypothesis $\mu^p=1$ then forces $\nu=1$.
\end{proof}

For a normalized MPS tensor, the transfer map is naturally trace-preserving but need
not be unital. In general one should not expect a single gauge choice to make it
both unital and trace-preserving.

\begin{remark}[Bi-canonical special case]\label{rem:bicanonical_special_case}\leanok
    When a Kraus map happens to be both unital and trace-preserving
    (the ``bi-canonical'' case), the Kadison--Schwarz equality at peripheral
    eigenvectors (Theorem~\ref{thm:ks_peripheral}) shows directly that the
    peripheral eigenvalues are closed under powers, and hence are roots of unity
    by Theorem~\ref{thm:peripheral_powers}. It
    requires both normalizations simultaneously. The more general adjoint-fixed-point
    formulation below covers the case where only unitality and a positive definite
    adjoint fixed point are available.
\end{remark}

\subsection{Peripheral closure via adjoint fixed point}%
\label{subsec:peripheral_closure_fixedpoint}

The Kraus map, its adjoint, and the unital normalization are those of
Definitions~\ref{def:kraus_map}, \ref{def:kraus_adjoint_map}, and
\ref{def:unital_kraus}. The following results combine the canonical
Kadison--Schwarz and multiplicative-domain theorems of Chapter~\ref{ch:schwarz}
with a positive definite fixed point of the adjoint map, replacing trace
preservation by the weighted-trace argument in
\cite[Appendix~A]{Cirac2016MPDO_arXiv}.

\begin{lemma}[Peripheral closure via adjoint fixed point]%
    \label{lem:peripheral_closed_powers_irred_fp}
    \lean{Kraus.peripheralEigenvalues_pow_mem_of_irreducible_unital_of_adjoint_fixedPoint}
    \leanok
    \uses{def:kraus_map, def:kraus_adjoint_map, def:unital_kraus,
        def:irreducible_cp}
    Let $E(X)=\sum_iK_iXK_i^\dagger$ be a Kraus map that is
    unital ($\sum_iK_iK_i^\dagger=\Id$), and assume:
    \begin{enumerate}
        \item the adjoint Kraus map $E^*(X)=\sum_iK_i^\dagger XK_i$
              has a positive definite fixed point $\rho>0$, and
        \item $E$ is irreducible.
    \end{enumerate}
    Then $\mathrm{peripheral}(E)$ is closed under powers:
    if $\mu\in\mathrm{peripheral}(E)$, then
    $\mu^n\in\mathrm{peripheral}(E)$ for every $n\in\N$.
\end{lemma}

\begin{proof}\leanok
    \uses{thm:kadison_schwarz, thm:psd_fp_pd_irred,
        thm:kraus_commute, thm:right_multiplicative_identity}
    Let $E(X)=\mu X$ with $|\mu|=1$, and set
    $G:=E(X^\dagger X)-E(X)^\dagger E(X)$.
    By (\ref{eq:schwarz_kadison}), $G\geq0$.
    Since $E^*(\rho)=\rho$,
    \begin{align}
        \tr(\rho G)
         & =\tr(\rho E(X^\dagger X))
        -\tr(\rho E(X)^\dagger E(X)) \notag \\
         & =\tr(E^*(\rho)X^\dagger X)
        -\tr(\rho E(X)^\dagger E(X)) \notag \\
         & =\tr(\rho X^\dagger X)
        -|\mu|^2\tr(\rho X^\dagger X)
        =0.
    \end{align}
    Because $\rho>0$ and $G\geq0$, this forces $G=0$, so
    $E(X^\dagger X)=E(X)^\dagger E(X)=X^\dagger X$.
    Hence $X^\dagger X$ is a positive semidefinite fixed point, which is
    positive definite by irreducibility
    (Theorem~\ref{thm:psd_fp_pd_irred}). Thus $X$ is invertible.
    Theorem~\ref{thm:kraus_commute} gives the corresponding
    intertwining relation with each Kraus operator. Iterating the
    multiplicative identity of
    Theorem~\ref{thm:right_multiplicative_identity} gives
    $E(X^n)=\mu^nX^n$, so $\mu^n$ is peripheral.
\end{proof}

\begin{theorem}[Roots of unity via adjoint fixed point]%
    \label{thm:peripheral_roots_irred_fp}
    \lean{Kraus.peripheral_isRootOfUnity_of_irreducible_unital_of_adjoint_fixedPoint}
    \leanok
    \uses{def:kraus_map, def:kraus_adjoint_map, def:unital_kraus,
        def:irreducible_cp, def:peripheral_eigenvalues}
    Let $E(X)=\sum_iK_iXK_i^\dagger$ be a Kraus map that is
    unital, and assume that the adjoint Kraus map
    $E^*(X)=\sum_iK_i^\dagger XK_i$ has a positive definite fixed
    point and that $E$ is irreducible.
    Then every peripheral eigenvalue of $E$ is a root of unity.
\end{theorem}

\begin{proof}\leanok
    \uses{lem:peripheral_closed_powers_irred_fp, thm:peripheral_powers}
    Combine the closure-under-powers result
    (Lemma~\ref{lem:peripheral_closed_powers_irred_fp})
    with Theorem~\ref{thm:peripheral_powers}.
\end{proof}

\begin{remark}\label{rem:adjoint_fp_more_general}\leanok
    Lemma~\ref{lem:peripheral_closed_powers_irred_fp}
    and Theorem~\ref{thm:peripheral_roots_irred_fp}
    cover the bi-canonical (unital + TP) case as well:
    trace preservation gives
    $E^*(\Id)=\Id$, which is a positive definite fixed
    point of the adjoint.
    The adjoint-fixed-point formulation matches the
    argument in \cite[Appendix~A]{Cirac2016MPDO_arXiv},
    which uses a faithful invariant state rather than
    bi-canonicality.
    Theorem~\ref{thm:peripheral_roots_irred_fp} establishes the
    root-of-unity conclusion in this adjoint-fixed-point formulation.
    In the same adjoint-fixed-point formulation, Chapter~\ref{ch:spectral}
    proves the corresponding cyclic structure: the cyclic group of
    peripheral eigenvalues is Theorem~\ref{thm:peripheral_cyclic_structure},
    and the cyclic projection decomposition is
    Theorem~\ref{thm:cyclic_decomposition_irreducible_schwarz}.
\end{remark}

\begin{definition}[Channel period]\label{def:channel_period}
    \lean{channelPeriod}
    \leanok
    \uses{def:peripheral_eigenvalues}
    The \emph{period} of a channel $E$ is the number of
    peripheral eigenvalues, counted without multiplicity:
    $p(E):=|\mathrm{peripheral}(E)|$.
\end{definition}

\section{Primitivity and the complementary transfer-map gap}

\begin{theorem}[Primitivity equals period one]\label{thm:primitive_iff_period_one}
    \lean{isPrimitive_iff_period_one}
    \leanok
    \uses{def:primitive_channel, def:channel_period}
    A linear map with a nonzero fixed point is primitive if and only if
    its period is~$1$.
\end{theorem}

\begin{proof}\leanok\uses{thm:one_peripheral}
    If $\mathrm{peripheral}(E)=\{1\}$, then
    $|\mathrm{peripheral}(E)|=1$.
    Conversely, if $|\mathrm{peripheral}(E)|=1$, then
    $1\in\mathrm{peripheral}(E)$ by
    Theorem~\ref{thm:one_peripheral}, so the only element is~$1$.
\end{proof}

\begin{remark}[Period-one characterization]%
    \label{rem:primitive_period_one_characterization}
    \leanok
    \uses{def:channel_period, thm:primitive_iff_period_one}
    The period $p(E)$ counts peripheral eigenvalues without multiplicity.
    Theorem~\ref{thm:primitive_iff_period_one} is therefore the period-one
    characterization of primitivity. The spectral estimates below are stated
    directly with Definition~\ref{def:primitive_channel}, since they use the
    absence of peripheral eigenvalues other than~$1$, not a separate counting
    hypothesis.
\end{remark}

\begin{theorem}[Complementary transfer-map gap for primitive maps]%
    \label{thm:primitive_compl_lt_one}
    \lean{compl_eigenvalue_norm_lt_one_of_primitive}
    \leanok
    \uses{def:primitive_channel, def:fixed_point_proj, def:trace_preserving}
    Let $E:\MN{D}\to\MN{D}$ be trace-preserving, and let
    $\rho\neq0$ satisfy $E(\rho)=\rho$ and $\tr(\rho)\neq0$.
    Let $P$ be the fixed-point projection associated to $\rho$, as defined
    in (\ref{eq:channel_fixed_point_proj}). Assume that:
    \begin{enumerate}
        \item $E$ is primitive;
        \item every eigenvalue $\mu$ of $E$ satisfies $|\mu|\le1$;
        \item whenever $E(X)=X$ and $\tr(X)=0$, one has $X=0$.
    \end{enumerate}
    Then every eigenvalue $\nu$ of $E-P$ satisfies $|\nu|<1$.
\end{theorem}

\begin{proof}
    \leanok
    Let $(E-P)(X)=\nu X$ with $X\neq0$. Then
    $E(X)=\nu X+P(X)$. If $\tr(X)=0$, then $P(X)=0$, so
    $E(X)=\nu X$ and hence $\nu$ is an eigenvalue of $E$ with
    $|\nu|\le1$. If $|\nu|=1$, primitivity forces $\nu=1$, so $X$ is a
    trace-zero fixed point of $E$; by assumption this implies $X=0$, a
    contradiction. Thus $|\nu|<1$.

    If instead $\tr(X)\neq0$, trace preservation and
    $\tr(P(X))=\tr(X)$, which follows from
    (\ref{eq:channel_fixed_point_proj}), give
    \begin{align}
        \tr(X)
         & =\tr(E(X))
        =\nu\tr(X)+\tr(P(X))
        =\nu\tr(X)+\tr(X).
        \notag
    \end{align}
    Hence $\nu\tr(X)=0$ and therefore $\nu=0$. In either case,
    $|\nu|<1$.
\end{proof}
